\documentclass[11pt]{article}
\usepackage[a4paper,margin=25mm]{geometry}
\usepackage{amsmath,amssymb,mathtools}
\usepackage{fontspec,unicode-math}
\setmainfont{Latin Modern Roman}
\setsansfont{Latin Modern Sans}
\setmonofont{Latin Modern Mono}
\setmathfont{Latin Modern Math}
\usepackage{xurl,hyperref}
\hypersetup{colorlinks=true,urlcolor=blue,linkcolor=blue}
\newcommand{\sref}[2]{\hyperlink{source:#1:#2}{[S#2]}}
\newcommand{\cataloguescope}[1]{\paragraph{Mathematical question.}#1}
\begin{document}
% UnsolvedMath ID 8700066; source catalogue code AMR-086-0066
\setcounter{section}{8700065}
\section{Conjecture 5.3}\label{8700066}
{\small\sffamily UnsolvedMath ID 8700066 · Graph Theory}
\subsection{Definitions and mathematical statement}

\paragraph{Shared notation.}$\mathbb N=\{1,2,\ldots\}$, and $\mathbb Z,\mathbb Q,\mathbb R,\mathbb C$ denote the integers, rationals, reals and complex numbers. Any other conventions are specified locally.


For an algebraic number $\alpha$, choose its primitive irreducible polynomial $f_\alpha\in\mathbb Z[T]$ with positive leading coefficient $a$ and roots $\alpha^{(1)},\ldots,\alpha^{(d)}$. Set
\[
 M(\alpha)=a\prod_{j=1}^d\max\{1,|\alpha^{(j)}|\},\qquad
 \mathrm h(\alpha)=\frac{\log M(\alpha)}d,
\]
and let $H(\alpha)$ be the maximum absolute value of the coefficients of $f_\alpha$. These are respectively Mahler measure, absolute logarithmic height, and usual polynomial height.

Almost all tuples here means outside a Lebesgue-null subset of $\mathbb R^n$, the real simultaneous-approximation setting of Section 5.1. Let $n\ge1$, $x=(x_1,\ldots,x_n)\in\mathbb R^n$, and write $F=\mathbb Q(\alpha_1,\ldots,\alpha_n,\beta_1,\ldots,\beta_n)$ for algebraic approximants.
\paragraph{Statement recorded in the supplied source.}
For almost every $x\in\mathbb R^n$, do there exist constants $c,D_0>0$, allowed to depend on $n$ and $x$, such that, for every integer $D\ge D_0$, every real $\mu\ge D$, and every $\alpha_1,\ldots,\alpha_n,\beta_1,\ldots,\beta_n\in\overline{\mathbb Q}$ with
\[
[F:\mathbb Q]\le D,\qquad [F:\mathbb Q]\max_{1\le i\le n}\{\mathrm h(\alpha_i),\mathrm h(\beta_i)\}\le\mu,
\]
one has
\[
\max_{1\le i\le n}\{|x_i-\beta_i|,|e^{x_i}-\alpha_i|\}\ge\exp(-cD^{1/(2n)}\mu)?
\] \sref{8700066}{2}

\subsection{Short English statement}
For almost every real tuple, is simultaneous algebraic approximation of the tuple and its exponentials uniformly limited by the stated field-degree and height bound?
\subsection{Sources}
\begin{description}
\item[\hypertarget{source:8700066:1}{S1}] ULAM.AI and the UnsolvedMath Contributors, UnsolvedMath: A Curated Collection of Open Mathematics Problems, record 8700066 (AMR-086-0066). CC BY 4.0. \url{https://huggingface.co/datasets/ulamai/UnsolvedMath}.
\item[\hypertarget{source:8700066:2}{S2}] Michel Waldschmidt, \emph{Open Diophantine Problems}, arXiv:math/0312440v2 (24 January 2004), Conjecture 5.3, p.47; real-tuple metric context, p.46; heights, p.35. \url{https://arxiv.org/abs/math/0312440v2}.
\end{description}
\label{end:8700066}
\end{document}
